\documentclass[12pt]{article}
\usepackage[utf8]{inputenc}
\usepackage[T1]{fontenc}
\usepackage{geometry}
\geometry{a4paper, margin=2.5cm}
\usepackage{setspace}
\onehalfspacing
\usepackage{hyperref}
\hypersetup{colorlinks=true, linkcolor=black, citecolor=black, urlcolor=black}
\usepackage{natbib}
\bibliographystyle{plainnat}
\usepackage{amsmath}
\usepackage{booktabs}

\title{The Kernel as an Open System:\\
\large Maxwell's Demon, Landauer's Bound, and the Lindblad Structure of a Running Viable System}
\author{[authors]}
\date{Skeleton — \today}

\begin{document}
\maketitle

\begin{abstract}
[TO WRITE — one paragraph per lens: (i) a working cybernetic kernel treated as a
thermodynamically open system; (ii) an empirical Maxwell demon on real dynamics
(ensemble cost $0.31$ bits vs.\ $500$ interventions, $1100\times$ discrepancy);
(iii) the Landauer accounting that restores the Second Law; (iv) the Lindblad
master-equation reading: the $\varphi^{\uparrow}/\varphi^{\downarrow}$ channel
pair as measurement-and-drive apparatus, the attractor as engineered
steady-state. Audience: complex systems / open quantum systems community
(Fierro UV, SFI).]
\end{abstract}

% ─── §1 ───────────────────────────────────────────────────────────────────────
\section{Introduction: Thermodynamic Information as a Common Language}
\label{sec:intro}
% MOVES: Jaynes' program (information thermodynamics) as the bridge language
% between cybernetics and statistical physics; why a RUNNING kernel (not a
% simulation of one) is the right object; the claim structure of the paper.
[SECTION TO WRITE]

% ─── §2 ───────────────────────────────────────────────────────────────────────
\section{The Shannon--Boltzmann Bridge}
\label{sec:bridge}
% MOVES: Jaynes 1957 formal equivalence; H(t) of the kernel state as physical
% entropy of an ensemble; what the equivalence licenses and what it does not.
[SECTION TO WRITE]

% ─── §3 ───────────────────────────────────────────────────────────────────────
\section{The VSF Kernel as an Open System}
\label{sec:kernel}
% MOVES: organelle decomposition; H(t) dynamics under load; phi-up as
% measurement channel / phi-down as drive; environment = host + adjacent LXCs;
% EXP-SOPLO-01 (selection entropy of the routing organ) slots here when run.
[SECTION TO WRITE]

% ─── §4 ───────────────────────────────────────────────────────────────────────
\section{Experiment: Maxwell's Demon on Real Dynamics}
\label{sec:demon}
% MOVES: EXP-OE-01 protocol + results (import from P1 §maxwell, DON'T repeat —
% cite and extend with the open-system reading); per-event vs ensemble cost.
[SECTION TO WRITE]

% ─── §5 ───────────────────────────────────────────────────────────────────────
\section{Landauer's Bound and the Restoration of the Second Law}
\label{sec:landauer}
% MOVES: where erasure actually happens in the kernel; the 1100x discrepancy
% resolved; cost concentrates at bifurcations (shared result with P_phi §2 —
% cite, different audience/register).
[SECTION TO WRITE]

% ─── §6 ───────────────────────────────────────────────────────────────────────
\section{The Lindblad Connection}
\label{sec:lindblad}
% MOVES: master equation dρ/dt = -i[H,ρ] + Σ_k (L_k ρ L_k† − ½{L_k†L_k, ρ});
% jump operators as algedonic/param_change events; the engineered steady-state
% AS the attractor; what the analogy predicts (spectral gap ↔ recovery rate)
% and how to falsify it on the kernel.
[SECTION TO WRITE]

% ─── §7 ───────────────────────────────────────────────────────────────────────
\section{The Quantum Computing Limit: Decoherence as the Unavoidable Demon}
\label{sec:quantum}
% MOVES: decoherence = environment measuring the register = a demon nobody
% hired; error correction as paying Landauer continuously; viability language
% for fault-tolerance thresholds.
[SECTION TO WRITE]

% ─── §8 ───────────────────────────────────────────────────────────────────────
\section{Discussion}
\label{sec:discussion}
[SECTION TO WRITE]

% ─── §9 ───────────────────────────────────────────────────────────────────────
\section{Open Questions and Future Experiments}
\label{sec:open}
% MOVES: EXP-SOPLO-01 (is the routing organ itself a demon?); thermal ramp
% recursion; Lindblad spectral-gap measurement on the real kernel.
[SECTION TO WRITE]

\bibliography{refs}
\end{document}
